\documentclass[11pt,a4paper]{article}

\usepackage[utf8]{inputenc}
\usepackage[T2A]{fontenc}
\usepackage[ukrainian,english]{babel}
\usepackage[top=22mm,bottom=22mm,left=22mm,right=22mm]{geometry}
\usepackage{hyperref}
\usepackage{fancyhdr}
\usepackage{parskip}
\usepackage{tcolorbox}
\usepackage{enumitem}
\usepackage{listings}
\usepackage{xcolor}

\tcbuselibrary{skins,breakable}

\hypersetup{
    pdftitle={Workshop 09 --- Agent SDK: Handout},
    pdfauthor={Vadym Bondarenko},
    colorlinks=true,
    linkcolor=blue,
    urlcolor=cyan,
}

\pagestyle{fancy}
\fancyhf{}
\fancyhead[L]{\small\textit{Workshop 09 --- Agent SDK}}
\fancyhead[R]{\small\thepage}
\fancyfoot[C]{\small\textit{vcdev.me \textperiodcentered{} 2026}}
\renewcommand{\headrulewidth}{0.4pt}

\lstdefinelanguage{python2}{
    keywords={import,from,as,def,return,if,else,elif,for,while,with,async,await,True,False,None,class,raise,try,except},
    sensitive=true,
    comment=[l]{\#},
    morestring=[b]',
    morestring=[b]"
}

\lstset{
    basicstyle=\ttfamily\small,
    breaklines=true,
    frame=single,
    backgroundcolor=\color{gray!10},
    rulecolor=\color{gray!40},
    xleftmargin=8pt,
    xrightmargin=8pt,
    aboveskip=6pt,
    belowskip=6pt,
}

\title{%
  {\Huge\bfseries Workshop 09: Agent SDK}\\[0.6em]
  {\large Будуємо власного асистента --- handout}
}
\author{Vadym Bondarenko \\ \small\href{mailto:vadym.bondarenko@vcdev.me}{vadym.bondarenko@vcdev.me}}
\date{2026}

\begin{document}

\maketitle

\begin{abstract}
Пост-воркшоп референс. Не повторює слайди --- сухий конспект SDK API, frontmatter-полів, code snippets, чек-листів. Кожен факт посилається на офіційні docs (\href{https://docs.claude.com/en/api/getting-started}{docs.claude.com}).
\end{abstract}

\tableofcontents
\newpage

\section{Три SDK для одного Claude}

\begin{tabular}{lll}
\hline
\textbf{Surface} & \textbf{Install} & \textbf{Use case} \\
\hline
Anthropic Client SDK & \texttt{pip install anthropic} & Кастомні цикли, прецизійний контроль \\
Claude Agent SDK & \texttt{pip install claude-agent-sdk} & Claude Code-like агенти \\
Claude Code CLI & \texttt{npm i -g @anthropic-ai/claude-code} & Інтерактивна робота \\
\hline
\end{tabular}

\textbf{Note:} Claude Code SDK перейменовано на Claude Agent SDK (\href{https://docs.claude.com/en/agent-sdk/overview}{docs}).

\textbf{Версія:} Opus 4.7 (\texttt{claude-opus-4-7}) потребує Agent SDK $\ge$ v0.2.111.

\section{Latest models (2026-04)}

\begin{tabular}{lllll}
\hline
\textbf{Model} & \textbf{API ID} & \textbf{Context} & \textbf{Input \$} & \textbf{Output \$} \\
\hline
Opus 4.7 & \texttt{claude-opus-4-7} & 1M & \$5/MTok & \$25/MTok \\
Sonnet 4.6 & \texttt{claude-sonnet-4-6} & 1M & \$3/MTok & \$15/MTok \\
Haiku 4.5 & \texttt{claude-haiku-4-5} & 200K & \$1/MTok & \$5/MTok \\
\hline
\end{tabular}

Source: \href{https://docs.claude.com/en/about-claude/models/overview}{models overview}.

\section{Anthropic SDK основа}

\subsection{Hello world}

\begin{lstlisting}[language=python2]
from anthropic import Anthropic

client = Anthropic()  # reads ANTHROPIC_API_KEY

message = client.messages.create(
    model="claude-opus-4-7",
    max_tokens=1024,
    messages=[{"role": "user", "content": "Hello"}],
)
print(message.content[0].text)
\end{lstlisting}

\subsection{Streaming}

\begin{lstlisting}[language=python2]
with client.messages.stream(
    model="claude-haiku-4-5",
    max_tokens=1024,
    messages=[{"role": "user", "content": "Write a haiku"}],
) as stream:
    for text in stream.text_stream:
        print(text, end="", flush=True)
    final = stream.get_final_message()
\end{lstlisting}

\subsection{Reading usage}

\begin{lstlisting}[language=python2]
u = response.usage
u.input_tokens                  # post-breakpoint input
u.output_tokens
u.cache_creation_input_tokens   # cache write (1.25x or 2x)
u.cache_read_input_tokens       # cache hit (0.1x)
\end{lstlisting}

\section{Tool use cheatsheet}

\subsection{Define}

\begin{lstlisting}[language=python2]
TOOLS = [{
    "name": "days_between",
    "description": "Calculate days between two ISO dates.",
    "input_schema": {
        "type": "object",
        "properties": {
            "start": {"type": "string"},
            "end":   {"type": "string"},
        },
        "required": ["start", "end"],
    },
}]
\end{lstlisting}

\subsection{Loop pattern}

\begin{lstlisting}[language=python2]
messages = [{"role": "user", "content": user_msg}]
MAX_ITER = 10
for i in range(MAX_ITER):
    response = client.messages.create(
        model="claude-opus-4-7", max_tokens=1024,
        tools=TOOLS, messages=messages,
    )
    if response.stop_reason != "tool_use":
        break
    tool_uses = [c for c in response.content if c.type == "tool_use"]
    tool_results = [
        {"type": "tool_result", "tool_use_id": tu.id,
         "content": [{"type": "text", "text": run_tool(tu.name, tu.input)}]}
        for tu in tool_uses
    ]
    messages.append({"role": "assistant", "content": response.content})
    messages.append({"role": "user", "content": tool_results})
\end{lstlisting}

\subsection{stop\_reason values}

\begin{itemize}
\item \texttt{end\_turn} --- normal completion, exit
\item \texttt{tool\_use} --- handle tool, loop
\item \texttt{max\_tokens} --- output truncated, raise limit
\item \texttt{stop\_sequence} --- stop sequence hit
\item \texttt{pause\_turn} --- server tool waiting, requery
\item \texttt{refusal} --- model refused, surface to user
\end{itemize}

\textbf{Pricing:} tools add 346 tokens (auto/none) or 313 tokens (any/tool) of system overhead per request for Opus 4.7, Sonnet 4.6, Haiku 4.5 (\href{https://docs.claude.com/en/agents-and-tools/tool-use/overview}{tool-use overview}).

\section{Prompt caching}

\subsection{Multipliers}

\begin{tabular}{lll}
\hline
\textbf{TTL} & \textbf{Write cost} & \textbf{Read cost} \\
\hline
5 min (default) & 1.25\,$\times$ input & 0.1\,$\times$ input \\
1 hour & 2.0\,$\times$ input & 0.1\,$\times$ input \\
\hline
\end{tabular}

\subsection{Syntax}

\begin{lstlisting}[language=python2]
response = client.messages.create(
    model="claude-opus-4-7", max_tokens=1024,
    system=[
        {"type": "text", "text": STATIC_PROMPT,
         "cache_control": {"type": "ephemeral"}},          # 5 min
        {"type": "text", "text": LARGE_DOC,
         "cache_control": {"type": "ephemeral", "ttl": "1h"}},
    ],
    messages=[{"role": "user", "content": question}],
)
\end{lstlisting}

\subsection{Limits}

\begin{itemize}
\item \textbf{Max breakpoints:} 4 per request
\item \textbf{Minimum cacheable:} 4096 tokens (Opus 4.7, 4.6, 4.5; Haiku 4.5); 2048 (Sonnet 4.6); 1024 (older)
\item \textbf{Lookback window:} 20 blocks per breakpoint
\end{itemize}

\subsection{Invalidation hierarchy}

\texttt{tools $\to$ system $\to$ messages}. Зміна вгорі $\to$ нижні рівні інвалідуються. Стратегія: статичне вгорі (tools, system, examples), динамічне внизу (last user message).

\subsection{Cost math example}

10-turn session, 50K-token system prompt, Opus 4.7 (\$5/MTok input):

\begin{itemize}
\item \textbf{Without caching:} $10 \times 50000 \times \$5/10^6 = \$2.50$
\item \textbf{With 5-min caching:} $50000 \times \$5 \times 1.25 / 10^6 + 9 \times 50000 \times \$5 \times 0.1/10^6 = \$0.31 + \$0.225 = \$0.535$
\item \textbf{Saved:} $\sim$78\%
\end{itemize}

\section{Claude Agent SDK}

\subsection{Hello world}

\begin{lstlisting}[language=python2]
import asyncio
from claude_agent_sdk import query, ClaudeAgentOptions

async def main():
    async for message in query(
        prompt="Find all TODO comments and summarize",
        options=ClaudeAgentOptions(allowed_tools=["Read", "Glob", "Grep"]),
    ):
        if hasattr(message, "result"):
            print(message.result)

asyncio.run(main())
\end{lstlisting}

\subsection{Built-in tools}

\texttt{Read}, \texttt{Write}, \texttt{Edit}, \texttt{Bash}, \texttt{Glob}, \texttt{Grep}, \texttt{Monitor}, \texttt{WebSearch}, \texttt{WebFetch}, \texttt{AskUserQuestion}.

\subsection{Two interfaces}

\begin{itemize}
\item \texttt{query()} --- one-shot, async iterator
\item \texttt{ClaudeSDKClient} --- multi-turn, custom in-process tools, hooks, sessions
\end{itemize}

\subsection{Filesystem config}

Agent SDK reads \texttt{.claude/} and \texttt{\textasciitilde/.claude/} by default: skills, commands, \texttt{CLAUDE.md}, plugins. Opt out with \texttt{setting\_sources=[]}.

\section{Debug top-6}

\begin{enumerate}
\item \textbf{\texttt{max\_tokens} forgotten} --- \texttt{TypeError} on \texttt{messages.create}
\item \textbf{\texttt{stop\_reason=max\_tokens}} --- output truncated, raise limit
\item \textbf{Cache not hitting} --- check size $\ge$ 4096, prefix unchanged, breakpoint position
\item \textbf{\texttt{asyncio} in Jupyter} --- \texttt{import nest\_asyncio; nest\_asyncio.apply()}
\item \textbf{\texttt{stop\_reason=refusal}} --- model refused, soften system prompt
\item \textbf{Tool loop infinite} --- enforce \texttt{MAX\_ITER}, log tool inputs
\end{enumerate}

\section{Production checklist}

\begin{itemize}[leftmargin=2em]
\item[$\square$] \texttt{MAX\_ITER} in tool loop ($\sim$10)
\item[$\square$] \texttt{max\_tokens} chosen consciously
\item[$\square$] Log \texttt{usage} on every call
\item[$\square$] \texttt{cache\_control} on static blocks $\ge$ 4096 tokens
\item[$\square$] Retry on \texttt{429} (SDK does this; verify config)
\item[$\square$] Timeout reviewed (default 600s)
\item[$\square$] Handle \texttt{RateLimitError}, \texttt{APIError}, \texttt{BadRequestError}
\item[$\square$] Handle all 6 \texttt{stop\_reason} values
\item[$\square$] API key from env / secret manager, never in code
\end{itemize}

\section{Setup commands}

\subsection{Install + auth}

\begin{lstlisting}[language=bash]
pip install anthropic claude-agent-sdk
export ANTHROPIC_API_KEY=sk-ant-...
\end{lstlisting}

\subsection{Workshop exercises}

\begin{lstlisting}[language=bash]
cd workshops/09-agent-sdk/exercises
python -m venv .venv && source .venv/bin/activate
pip install -r requirements.txt
python solutions/01-hello-agent/hello.py
python solutions/02-tool-use/agent.py
python solutions/03-prompt-caching/cache_demo.py
python solutions/04-mychat-cli/mychat.py
\end{lstlisting}

\section{Resources}

\begin{itemize}
\item \href{https://docs.claude.com/en/api/getting-started}{docs.claude.com/en/api/getting-started}
\item \href{https://docs.claude.com/en/agent-sdk/overview}{docs.claude.com/en/agent-sdk/overview}
\item \href{https://docs.claude.com/en/agents-and-tools/tool-use/overview}{docs.claude.com/en/agents-and-tools/tool-use/overview}
\item \href{https://docs.claude.com/en/build-with-claude/prompt-caching}{docs.claude.com/en/build-with-claude/prompt-caching}
\item \href{https://docs.claude.com/en/about-claude/models/overview}{docs.claude.com/en/about-claude/models/overview}
\item \href{https://github.com/anthropics/anthropic-sdk-python}{github.com/anthropics/anthropic-sdk-python}
\item \href{https://github.com/anthropics/claude-agent-sdk-python}{github.com/anthropics/claude-agent-sdk-python}
\item \href{https://github.com/anthropics/claude-agent-sdk-demos}{github.com/anthropics/claude-agent-sdk-demos}
\end{itemize}

\end{document}
